\documentclass[a5paper,10pt]{book}

\usepackage{amsfonts}
\usepackage{hyperref}
\usepackage[utf8]{inputenc}
\usepackage[dvipsnames]{xcolor}
\usepackage[hmarginratio=1:1]{geometry}% equal left and right margins
\usepackage{graphicx}
% \usepackage{fontspec}
% \setmainfont{Times New Roman}
\usepackage{titling}

% \usepackage{sectsty}
% \chapterfont{\color{Red} \normalsize \mdseries}  % sets colour of chapters
% \sectionfont{\color{Red}}  % sets colour of sections

\usepackage{afterpage}
\newcommand\myemptypage{
    \null
    \thispagestyle{empty}
    \addtocounter{page}{-1}
    \newpage
    }

\setlength{\droptitle}{-1in}

\def\ae{æ}
\def\AE{Æ}


\begin{document}
\pagenumbering{roman}


\title{\includegraphics[width=2in]{Picture1.png} \\ \Huge \bfseries \color{Red} B R E V I A R V M \\ \huge R O M A N V M, \\ \large \mdseries \color{Black} Ex sacra potissimum scriptura, \& probatis sanctorum \\ historiis nuper consectum, ac denuo per \\ eundem authorem accurantius \\ recognitum. \\ \ \\ \itshape scrutamini scripturas, quoniam illæ sunt, \\ quæ testimonium perhibent de me. \\Iohan. V.}

\author{VENETIIS APVD IVNCTAS}
\date{MDCXXXI}

\maketitle

\myemptypage

% \frontmatter

\pagenumbering{roman}

\begin{center} \color{red}
AD SANCTISSIMVM PATREM \\ \& Dominum nostrum Paulum tertium \\ Pontificem Maximum, \\ \color{black} Francisci Quignonij Tituli Sanctæ Crucis in Ierusalem presbyteri Cardinalis, in Breuiarium proxime consectum, ac denuo \\ recognitum.
\end{center}

\centerline{\color{red}P R \AE \ F A T I O.}

\Huge \textbf{B}
\normalsize Reviarum Romanum nuper a nobis felicis recordationis Clementis. VII. Pontifi. Maximi hortatu consectum, ac potius in ampliorem sacrarum scripturam lectionem, ad veterem sanctorum patrum, \& conciliorum antiquorum formam reuocatum, tuaque voluntate sanctissime pater editum, graues plerosque ac doctos viros ita probasse \& recepisse intellexi, vt nihil in eo metandum existimarent.
Alios item animaduerti graues etiam, \& prudentes homines, qui eius rationem magnopere probantes, nonnihil tamen in eo desiderari affimarent.
Illud vero nunquam dubitatui, fore in tanta multitudine nonnullos ex ijs videlicet, qui in diuerso precandi ritu consenuissent, quibus labor ille noster non esset perinde gratus, existimantibus ab inueterata illa consuetudine precandi nulla ratione clericis esse discedendum.
Imo vero nobis primam editionem Breuiarij, non tanquam promulgationem legis esse placuerat, sed quasi publicam quandam deliberationem:
vt sic proposita nostra sententia, iudicia multorum exquireremus, \& quod omnium commodissimum, \& religioni, ac pietati conuenientissimum plerisque prudentibus, grauibusque viris visum esset, sequeremur.
Nihil enim humano elaboratum ingenio tam exactum initio vnquam fuit, quin postea multorum accedente iudicio perfectius reddi possit, vt in ipsis etiam ecclesiasticis institutis, circa primitiuam presertim ecclesiam contigisse videmus.
Itaque, multorum sententiis % end pp7
collatis, qu\ae \ nobis partim vocibus, partim scriptis innotuerunt, iudicium eorum secuti, qui omnium prudentissime sentire visi sunt, libenter quedam addidimus, alia mutauimus, \& omnia dilidenter recongnouimus, retenta tamen summa forma Breuiarii.
Sed quoniam sic fert natura rerum, vt nihil sit tam rectum, nihil ranta ratione in vitam vsumque hominum inductum, cuius nouitas non sit aliquibus ingrata, non temere facturi esse videmur, si rationem totius instituti nostri a nobis prius summatim redditam, nunc accuratius recognito Breuiario, paulo latius explicabimus.
Mihi enim (vt s\ae pe sum professus) cogitanti, atque animo repetenti initia veteris insituti, quo sancitum est, vt clerici sacris initiati, vel sacerdotiis pr\ae sidentes, singulis diebus perlegant horarias preces, quas Canonicas etiam appellamus, tres pr\ae cipue caus\ae \ spectat\ae \ fuisse videri solent.
Quarum ea prima est, qui cum c\ae teri homines in quaque ciuitate, aut suum quisque negocium agant, aut in Republica administranda sint occupati, clericis ex eo vocatis, qui sortis Domini sint, quique bonis ecclesiasticis aluntur, hoc potissimum negocium diuinis, \& humanis legibus est iniunctum, vt Deum habere propitium in comissum sibi populum, de se bene merentem cunctis rationibus enitantur.
quod %sic
non solum sacrificiis efficitur, sed etiam precibus, qu\ae \ a pio corede proficiscantur.
ipsi %sic
enim (vt testimonio versibusque vtar beati Alexandri pap\ae \ \& martyris) pro populo interpellant, \& populi peccata comedunt, quia precibus suis \& oblationibus ea delent, atque consumunt.
Qui quanto digniores fuerint, tanto facilius pro quibus necessitatibus clamant, exaudiuntur.
Quod testimonio probat Iacobi Apostoli, qui nos ad precandum cohortans, Orate (inquit) pro inuicem vt saluemini.
multum %sic
enim valet deprecatio iusti assidua.
Altera causa est, vt qui reliquo populo exemplo debent esse virtutis \& sanctimoni\ae , assidua precatione Deum alloquentes, minus opportuni reddantur tentatori diabolo, si eos inuenerit, vt Hieronumus ait, occupatos, \& a cogitationibus caducarum rerum subinde auocati, contemplationi diuinarum assuescant.
Tertia, vt religionis quoque futuri magistri quotidiana sacr\ae \ scriptur\ae \ \& ecclesiasticum historiarum lectione erudiantur, complectanturque (vt Paulos ait) eum, qui secundum doctrinam est, fidelem %end pp8
sermonem, \& potentes sint exhortari in doctrina sana, \& eos, qui contradicunt, arguere.
Et profecto si quis modum precandi olim a maioribus traditum diligenter confideret, plane inelligat horum omnium pr\ae cipia ab ipsis habitam esse rationem.
Sed factum est nescio quo pacto precantium negligentia, vt paulatim a sanctissimis illis veterum patrum institutis discederetur.
Nam libri scripure sacr\ae, qui statis anni temporibus legendi erant more maiorum (vt est in monumentis Gelasii pap\ae, \& concilii Romani. lxx. episcoporum) vixdum incepti omittuntur in alio Breuiario.
Tam histori\ae \ sanctorum quedam tam inculte, \& tam sine delectu script\ae \ habentur in eodem, vt nec authoritatem habere videantur, nec grauitatem.
Accedit tam perplexus ordo, tamque difficilis precandi ratio, vt interdum paulo minor opera in requirendo ponatur, qui, cum inueneris, in legendo.
Quibus rebus animaduersis, felicis recordationis Clemens. VII. Pontifex Maxi. cum intelligeret officii sui esse, cum aliorum Christianorum commoditatibus prospicere, tum in primis clericorum, quibus ministris vteretur in commisso sibi grege administrando, me hortatus est, negociumque dedit, vt quantum cura, \& diligentia niti possem, preces horarias ita disponerem, \& ad veterm illam formam, quoad fieri posset, redigerem:
vt difficultatibusm dispendiisque sublatis, nec a summa rationis olim ab antiquis sanctisque patribus institut\ae, discederetur:
nec rursum clerici ab officio precandi deterrerentur laboris magnitudine.
Quam ego pro uinciam libentissime suscepi, simul vt Christi vicario honestissima precipienti parerem, simul vt bono publico pro mea virili parte seruirem.
Adhibitis igitur quibusdam meorum domesticorum prudentibus hominibus, ac sacrarum literarum \& pontificii iuris doctrina pr\ae ditis, eisdemque Gr\ae ce \& Latine eruditis, dedi operam qui maxime potui, vt commode ac ex vtilitate publica rem conficerem.
Ac illud ante omnia visum nobis est in consuetudinem reuocare, vt scriptura sacra 



\newpage
\begin{center} \color{red}
Tabula literarum Dominicalium. 1564.
\end{center}

\begin{tabular}{c c c c c c c c c c c c c c}
b & g & f & e & d & b & \color{red} A & g & f & d & c & b & \color{red} A & f\\
\color{red} A & \  & \  & \  & c & \  & \  & \  & e & \  & \  & \  & g & \ \\
e & d & c & \color{red} A & g & f & e & c & b & \color{red} A & g & e & d & c\\
\  & \  & b & \  & \ & \  & d & \  & \  & \  & f & \  & \  & \ \\
\end{tabular}

\Huge \textbf{\color{red} I}
\normalsize N Tabula suprascripta, literam Dominicalem singulis annis hoc pacto inuenius.
Primam literam, qu\ae \ est \color{red} b a \color{black} tribue anno. \color{red} 1564.
\color{black} Sequentem, qu\ae \ est \color{red} g \color{black} tribue anno. \color{red} 1565.
\color{black} Sicque deinceps progredere quousque peruenias ad numerm anni, cuius literam scire vis.
Quia litera super quam cadit numerus illius anni, erit litera Dominicalis.
Qu\ae \ si vnica occurrerit, annum communem esse intelliges:
si duplex, Bissextiliem.
Et tunc prior, id est superior, usque ad festim Mathi\ae \ apostoli feseruiet.
Inferior autem ad reliquam nni partem accommodabitur.
Sciendum autem qui finitis \color{red} 18. \color{black} annis, redeundum est ad principium cum annis tunc currentibus.
\begin{center} \color{red}
Aureus numeris.
\end{center}


\chapter{Fundamentals }
\vspace{-5em} 

\section{Vectors in Euclidean Space}

% \mainmatter

\pagenumbering{arabic}

\chapter{First Chapter}

This will be an empty chapter...

\newpage
\thispagestyle{empty}
\mbox{}
\newpage

\chapter{Fundamentals of Vectors}

\section{Vectors in Euclidean Space}







\end{document}

